\chapter{Methodology}
\label{ch:methodology}

\section{Overview of the research design}
\label{sec:overview}

This thesis asks a simple question. When capable and limited language models share an
economy, who ends up better off, and does the ability to talk change the answer? To study
it, we place large language model (LLM) agents of two capability tiers into iterated,
mixed-motive economic games and treat their behavior as the measured outcome: what they
contribute, what they say, and what they keep. Capability is the variable we manipulate;
everything else is what we measure.

The design is factorial and fully pre-registered. We cross three factors: the composition
of the population (how many strong and weak agents meet in a game), the game itself (a
four-player Public Goods Game and a two-player Trust Game), and communication (whether
agents may post a public message before they act). These factors produce sixteen
experimental cells, each replicated over twenty random seeds, for a corpus of 320 games.
The strong tier is not one model but three independently developed frontier models, Claude
Haiku~4.5, GPT-5-mini, and Gemini~2.5~Flash, rotated through the strong slots so that no
result rests on the quirks of a single lab. The weak tier is a small model run locally,
Llama~3.2~3B.

Two choices carry most of the methodological weight. First, every agent runs at identical
inference settings, so a behavioral gap between tiers reflects capability and not a
difference in sampling temperature or decoding. Second, the study was registered on the
Open Science Framework before the first production game ran: the hypotheses, the planned
sample, the analysis, and the criteria that would falsify each claim. That pre-registration
is the backbone of the work. It commits us in advance to bounded predictions, including the
ones that later proved modest or null, and it is what lets the results chapter report a
sharp effect in one game and its near-absence in another without either looking like a
story chosen after the fact.

From each game we take four families of measurement. Welfare, the payoffs agents
accumulate, gives the inequality indices and the capability premium. Messages, the text
agents exchange, feed a message classifier and a contemporaneous-surprisal measure of
deception. Revealed actions feed a classifier that tries to recover an agent's tier from
behavior alone. Transfers between agents, in the games that have them, form an
interaction network. The rest of the
chapter specifies each in turn: the experimental factors, the capability tiers and their
rotation, the game protocols, the inference settings, the metrics, the sample and budget,
the pre-registration and analysis plan, and the integrity safeguards and limitations of
the design.

\section{Experimental factors: composition, game, communication}
\label{sec:factors}

Three factors define the experiment. A single game is one combination of a population
composition, a game type, and a communication condition. Crossing the three produces the
sixteen cells in Table~\ref{tab:matrix}, each replicated over twenty seeds.

\textbf{Composition} sets how many strong and weak agents share a game. The homogeneous
compositions set a baseline for what each tier does among its own kind: all-strong and
all-weak in the Public Goods Game, strong-strong and weak-weak in Trust. The symmetric
mixes (two strong and two weak, or one of each in Trust) put the tiers on equal numerical
footing. The asymmetric compositions push the extremes: \emph{lone-elite} drops a single
strong agent into a group of weak ones, and \emph{lone-vulnerable} does the reverse, a
single weak agent among the strong. These two ask whether a lone capable agent can dominate
a weak crowd, and whether a lone weak agent is exploited by the capable ones.

\textbf{Game} is either the Public Goods Game, played by four agents, or the Trust Game,
played by two. The choice is deliberate. The Public Goods Game is an $n$-player contribution
to a shared pool, where the temptation is to free-ride on everyone else; the Trust Game is
a two-player exchange that turns on whether one agent will repay another's trust. Running
both lets us ask not just whether capability matters, but whether an effect belongs to
group dynamics or to one-on-one trust, a distinction that turns out to matter
(Chapter~\ref{ch:results}). Both games run for ten rounds.

\textbf{Communication} is either off or on. When it is on, each agent may write one
short public message before choosing its action that round, framed to it as visible to
the others. The message is cheap talk: non-binding and costless. In this build it is
produced and logged but not delivered to the other agents (Section~\ref{sec:communication}),
so the condition is a public commitment without reception rather than an exchange. It is
the single lever the headline result turns on, and switching it is how the study isolates
the effect of asking an agent to commit in public.

\begin{table}[t]
\centering
\caption{The sixteen-cell experimental matrix. Each population composition is played in one
of the two games under both communication conditions, with twenty seeded replications per
cell, for 320 games in total.}
\label{tab:matrix}
\begin{tabular}{llccc}
\toprule
Game & Composition & Strong\,:\,Weak & Comm.\ off & Comm.\ on \\
\midrule
\multirow{5}{*}{Public Goods ($n=4$)}
  & all-strong      & 4\,:\,0 & 20 & 20 \\
  & all-weak        & 0\,:\,4 & 20 & 20 \\
  & mixed           & 2\,:\,2 & 20 & 20 \\
  & lone-elite      & 1\,:\,3 & 20 & 20 \\
  & lone-vulnerable & 3\,:\,1 & 20 & 20 \\
\midrule
\multirow{3}{*}{Trust ($n=2$)}
  & strong-strong   & 2\,:\,0 & 20 & 20 \\
  & strong-weak     & 1\,:\,1 & 20 & 20 \\
  & weak-weak       & 0\,:\,2 & 20 & 20 \\
\midrule
\multicolumn{3}{l}{\textbf{Total}} & \multicolumn{2}{c}{\textbf{320} games} \\
\bottomrule
\end{tabular}
\source{Author}
\end{table}

\section{Capability tiers and model rotation}
\label{sec:tiers}

The two tiers are defined by capability, not by vendor. The weak tier is a single small
model, Llama~3.2~3B (\texttt{llama3.2:3b}), run locally through Ollama at 4-bit
quantization on a 4\,GB consumer GPU. The hardware set this choice: such a card runs a 3B
model at usable speed and nothing much larger, which is a fair picture of the weak agent a
budget-constrained deployment would actually field. The strong tier is three frontier
models from three different labs: Claude Haiku~4.5 (\texttt{claude-haiku-4-5-20251001}),
GPT-5-mini (\texttt{gpt-5-mini}), and Gemini~2.5~Flash (\texttt{gemini-2.5-flash}).

The strong tier is three models rather than one because of a flaw in the pilot. An early
pilot used Haiku alone as the strong agent and produced clean results, but a result from a
single model is a result about that model. Anything the pilot credited to ``the strong
tier'' could equally have been a quirk of Haiku's training, its refusal style, or its
preferred output format. This is pseudoreplication: many games, but one underlying strong
agent beneath all of them.

Production removes that confound by rotating the three strong models through the strong
slots, game by game, so the strong positions across the corpus are shared among
independently built systems. The rotation is balanced. Haiku fills 178 strong slots, Gemini
174, and GPT-5-mini 168, of 520 strong positions in total. Because the three models come
from different labs trained on different data, an effect that holds in all three is
hard to pin on any single model. When the results chapter reports that
communication flips the capability premium, the claim rests on three models moving the same
way, not on one model moving while two are assumed to follow.

Rotation also surfaces something a single-model design would hide: the strong tier is not
monolithic. The three models do not play alike. GPT-5-mini, for one, contributes less than
Haiku or Gemini in the same opening rounds. Treating the tier as a rotation rather than a
single model keeps that spread visible, and keeps the strong-tier claims honest about their
own internal variation.

\section{Game protocols and parameters}
\label{sec:protocols}

Both games run for ten rounds, and in both, holding back is privately tempting while
contributing is collectively better. The parameters are collected in
Table~\ref{tab:params}.

\begin{table}[t]
\centering
\caption{Game parameters. Both games are iterated for ten rounds; the per-round structure
is what differs.}
\label{tab:params}
\begin{tabular}{lll}
\toprule
Parameter & Public Goods Game & Trust Game \\
\midrule
Players & 4 & 2 (investor, trustee) \\
Rounds & 10 & 10 \\
Endowment / round & 20 tokens & 10 tokens (investor) \\
Multiplier & 1.6, on the pooled total & 3, on the amount sent \\
Action & contribute $c \in [0, 20]$ & send $x \in [0, 10]$; return $y \in [0, 3x]$ \\
Roles & symmetric & investor / trustee, alternating \\
Communication & off or on (cheap talk) & off or on (cheap talk) \\
\bottomrule
\end{tabular}
\source{Author}
\end{table}

In the \textbf{Public Goods Game}, four agents play together. Each round, every agent
receives an endowment of twenty tokens and chooses how much of it to contribute to a common
pool, keeping the rest. The pool is multiplied by 1.6 and split equally among all four, so
an agent's payoff for the round is the tokens it kept plus its quarter-share of the
multiplied pool, where the sum runs over all four players including agent~$i$ itself:
\[
  \pi_i = (20 - c_i) + \frac{1.6 \sum_{j} c_j}{4}.
\]
Because the multiplier divided by the group size is below one ($1.6/4 = 0.4$), each token an
agent contributes returns less than a token to itself, yet every contribution grows the
pool for the group. That gap is the temptation to free-ride.

In the \textbf{Trust Game}, two agents play, one as investor and one as trustee. The
investor receives ten tokens and sends some amount $x$ to the trustee. The amount triples in
transit, so the trustee receives $3x$ and returns some amount $y$. The investor ends the
round with $10 - x + y$ and the trustee with $3x - y$. Sending is a risk for the investor
and repaying is a cost to the trustee, so trust pays off only when the trustee chooses to
honor it. Roles alternate each round, so over ten rounds each agent plays investor five
times and trustee five times.

Each round, an agent sees the anonymized history of play so far and returns its action as a
structured (JSON) object: a contribution in the Public Goods Game, or an amount to send or
return in the Trust Game. When communication is on, the object also carries one short public
message, which the prompt frames as visible to the others, though the runner logs it
without delivering it (Section~\ref{sec:communication}). Agent identifiers are
anonymous eight-character strings, and the order in which agents appear in a prompt is
randomized every round, so no agent can be followed across games by name or treated as
higher-status because of where it sits in the list. If a response fails to parse, the agent
is re-prompted; after repeated failures it falls back to a default of contributing, sending,
or returning nothing.

\section{Inference settings (uniform, to isolate capability)}
\label{sec:inference}

The study compares capability, so the inference settings must not become a second, hidden
variable. If one model sampled at a different temperature, or was allowed a longer private
reasoning budget, a behavioral difference could reflect that setting rather than the model.
Every agent therefore runs at the same settings, and where one model forced a choice, the
others were brought into line with it rather than the reverse.

Temperature is fixed at 1 for every agent. This was not the original default; it was forced
by GPT-5-mini, which accepts only temperature 1 and returns an error at lower values. Rather
than leave GPT-5-mini at 1 while the others sampled lower, which would be a configuration
confound, all agents were set to temperature 1, so that an observed gap reflects capability
and not a difference in sampling. The choice is recorded in the pre-registration.

Two of the strong models can spend tokens on a private reasoning pass before they answer,
and both were configured not to. GPT-5-mini runs at reasoning effort \texttt{minimal}, and
Gemini at a thinking budget of zero. There were two reasons. The practical one: at its
default effort GPT-5-mini spent roughly nine hundred reasoning tokens before producing any
output, overran the 512-token response budget, and returned nothing, which failed to parse
on every call of an early smoke test. The principled one: a hidden chain of reasoning would
make the comparison unfair, because the other models answer directly. Holding all of them
to a single visible response keeps the comparison about the answer each model gives, not
about how much private computation it was granted. The weak model, served locally through
Ollama, has no such setting; it is steered toward valid output by requesting JSON at the
prompt level and re-prompting on a malformed response, up to three attempts before the
default action is used.

One substitution is worth recording. The intended Gemini checkpoint
(\texttt{gemini-3.5-flash}) was unavailable at launch, returning capacity errors on most
calls, while \texttt{gemini-2.5-flash} served reliably throughout the dry run. Production
therefore used \texttt{gemini-2.5-flash}, a substitution declared in advance in the
pre-registration, with the rule that if it too became unavailable mid-run the current Gemini
Flash endpoint would be used and its exact string logged. Because each game is isolated,
only the affected Gemini games would have needed re-running; none did.

\section{Metrics}
\label{sec:metrics}

\subsection{Welfare inequality (Gini, Atkinson) and the capability premium}
\label{subsec:inequality}

Inequality and the capability premium both read off a single quantity: what each
agent walks away with. Writing $\pi_{i,t}$ for the payoff of agent~$i$ in
round~$t$ (Section~\ref{sec:protocols}, with the round index now made explicit),
an agent's cumulative payoff over a game is
\[
  \Pi_i = \sum_{t=1}^{10} \pi_{i,t}.
\]
Within one game this gives a short welfare vector, one entry per agent: four
entries in the Public Goods Game, two in the Trust Game. Every measure in this
subsection is computed on that vector, per game, and only then averaged across
the twenty seeded games of a cell, one of the sixteen conditions of
Table~\ref{tab:matrix}. Keeping the unit of analysis at the whole
game, rather than pooling every agent from every game into one distribution, is
what makes the confidence intervals defined below interpretable.

To measure how unevenly a game's winnings are spread we use the Gini
coefficient. On a welfare vector of $n$ agents with mean cumulative payoff
$\bar{\Pi}$ it is
\[
  G = \frac{\sum_{i}\sum_{j} |\Pi_i - \Pi_j|}{2\,n^2\,\bar{\Pi}},
\]
the average absolute gap between every pair of agents, scaled to lie in $[0,1]$.
A value of zero means every agent ended with the same amount; the index rises
toward $(n-1)/n$ as one agent captures the entire pie. We use the population form,
dividing by $n^2$ with no small-sample correction: with only four or two agents
in a game any correction would be a judgment call, and the per-game value is
averaged across the cell's twenty seeded games before we read anything into it.
That ceiling, $(n-1)/n$, is $0.75$ in the four-agent Public Goods Game and $0.5$
in the two-agent Trust Game, so Gini magnitudes are read within a game type,
never compared across the two.

Alongside the Gini we compute the Atkinson index at three inequality-aversion
settings, $\varepsilon \in \{0.5, 1, 2\}$,
\[
  A_\varepsilon = 1 - \frac{1}{\bar{\Pi}}
  \left( \frac{1}{n}\sum_{i} \Pi_i^{\,1-\varepsilon} \right)^{1/(1-\varepsilon)},
\]
with the $\varepsilon = 1$ case taking the limit form, where the bracket becomes
the geometric mean $\exp\big(n^{-1}\sum_i \ln \Pi_i\big)$. A higher $\varepsilon$
weights the bottom of the distribution more heavily, so the three settings
together test whether an inequality ordering survives being re-weighted toward the
worst-off. The Atkinson index was pre-registered alongside the Gini as a primary
inequality measure; Chapter~\ref{ch:results} reports it as a robustness check on
the Gini ranking.

In the two-player Trust Game the inequality indices collapse to a single unsigned
gap between two agents, so there the comparison leans instead on the
\emph{capability premium}, the difference in mean cumulative payoff between the two
tiers within a game:
\[
  \Delta = \bar{\Pi}_{\mathrm{strong}} - \bar{\Pi}_{\mathrm{weak}}.
\]
Unlike the Gini, it keeps the sign of that gap, which separates exploiter from
exploited. The premium is the main mixed-tier measure in both games, defined
wherever both tiers are present; the homogeneous baselines have no premium to
report. A positive $\Delta$ means the strong agents out-earned the weak ones, a
negative $\Delta$ that the weak agents did better. The negative case is the
inversion that Chapter~\ref{ch:results} examines. We report $\Delta$ in raw tokens
rather than as a normalized fraction. Within a game type the endowment is fixed,
so a token is a token and the differences compare cleanly; normalizing would only
matter for setting the two game types on a common scale, which we instead avoid by
never averaging a premium across game types.

A cell holds twenty seeded games, so each Gini and each premium is a mean over
twenty per-game values. We attach a 95\% confidence interval to that mean with a
percentile bootstrap of ten thousand resamples. The resampling unit is the game:
we draw games with replacement from the cell and recompute the mean, so the
interval reflects how far the statistic travels from one seed to the next, not
the spread among agents inside a single game. The bootstrap seed is fixed, so the
intervals are reproducible from the configuration alone. These are the intervals
drawn as whiskers on the inequality and premium figures in
Chapter~\ref{ch:results}.

\subsection{Deception signals: contemporaneous surprisal vs.\ global policy-KL}
\label{subsec:deception}

Deception here is a gap between words and deeds: an agent that promises to
cooperate and then free-rides. Putting a number on
that gap is harder than it sounds, because the obvious measure can track the
wrong thing. We measure it with a contemporaneous surprisal signal, a global
policy divergence kept as a control, and a blind text classifier as an
independent check. Both numeric measures were tested against human labels fixed
before any data existed, under a threshold set in advance; reporting the one that
failed beside the one we kept is the comparison this thesis treats as a result in
its own right.

The primary signal is \emph{contemporaneous surprisal}. In a communication round
agent~$i$ first sends a message, then chooses an action~$a_{i,t}$. A separate
scoring model (Claude Haiku~4.5) reads that message and infers the distribution
$q_{i,t}$ over the action the message implies, the agent's \emph{stated policy}
for the round. Surprisal is the negative log probability that this stated policy
assigned to the action actually taken,
\[
  s_{i,t} = -\ln q_{i,t}(a_{i,t}),
\]
averaged over the agent's message-bearing rounds. It is small when the action
lands where the words said it would, and large when the agent did something its
own message gave little weight to, the mark of saying one thing and doing
another. The actions live on a discrete grid (steps of five tokens in the Public
Goods Game, two in the Trust Game), and the stated distribution is given a small
floor so the logarithm stays finite. The measure is \emph{contemporaneous}
because each message is scored against the same round's action, not against the
agent's behavior pooled over the whole game.

That pooling is what the second measure does, and the reason we keep it only as a
control. \emph{Global policy-KL} compares the agent's revealed action
distribution over the entire game, $p_i$, against the same per-round stated
policies, through $D_{\mathrm{KL}}(p_i \,\|\, q_{i,t})$ averaged over the message
rounds. The trouble is that this rewards a peaked action distribution with a high
score: an agent that reliably plays the same cooperative action has a sharp
$p_i$, and against the spread-out distribution a hedged message implies, that
sharpness reads as large divergence. The most consistent cooperators therefore
score as the most deceptive, which is backwards. Global policy-KL measures how
concentrated a policy is, not whether words match deeds, so we register it as a
\emph{negative control}: it shows that not just any divergence between stated and
revealed behavior captures deception, only the contemporaneous, action-by-action
kind.

Keeping a measure we expect to fail is deliberate, and it follows from how the
study was pre-registered. Global policy-KL was the original confirmatory
deception metric, with a falsifiable validation threshold, a Spearman correlation
of at least $0.40$ against human labels, frozen before any validation. In a pilot
we hand-labeled one hundred messages blind, on a fixed zero-to-three deception
rubric, and correlated the human scores against both measures. Global policy-KL
showed no association ($\rho \approx 0$) and failed; contemporaneous surprisal
correlated with the human labels weakly but significantly ($\rho \approx 0.22$,
$p \approx 0.03$), yet still below the frozen $0.40$ bar.
Because surprisal cleared significance but not the confirmatory threshold, we
registered it as \emph{exploratory} rather than claim more than the validation
supported, and demoted policy-KL to the negative control described above. The
single human rater was the metric's author, mitigated by blind, single-pass
labeling against the fixed rubric and recorded as a limit rather than hidden.
Reporting the threshold that even surprisal fell short of, rather than quietly
dropping it, is part of that disclosure.

A third signal corroborates the first without sharing its machinery. A blind
message classifier reads each message and labels it, but never sees the action
behind it, so it cannot inherit any error from the action-based surprisal. Its
category for misleading statements gives an independent, text-only count of
deceptive messages by tier. The full label scheme and its distribution belong to
the message-classification metric (Section~\ref{subsec:classification}); here it
matters only as a third, methodologically independent check, so the deception
finding can rest on agreement between measures rather than on any one of them.
The production values for all three, and the transcripts behind them, are
reported in Chapter~\ref{ch:results}.

\subsection{Message classification}
\label{subsec:classification}

Where the deception signals of Section~\ref{subsec:deception} ask whether a
message was honest, this metric asks what kind of message it was. A classifier
(Claude Haiku~4.5, at temperature zero for stable labels) reads each message
together with its game context, never the action behind it, and sorts it into
one of five communicative types:
\begin{description}
  \item[cooperative-promise:] a first-person commitment to act pro-socially;
  \item[threat:] a warning of punishment contingent on what others do;
  \item[informational:] strategy or facts about the game, with no promise or
    threat;
  \item[deceptive-claim:] a statement about one's own intentions or the game that
    reads as misleading or manipulative;
  \item[neutral:] greeting or filler with no strategic import.
\end{description}
The output is a distribution over these five types for each tier, a picture of how
the strong and weak agents talk.

Because the classifier reads only the words, it shares no input with the
action-based surprisal of Section~\ref{subsec:deception}, so where the two agree
the corroboration is independent. What it adds is a division of labor: a message's
class is a claim about its surface form, not about whether the agent later honored
it, and whether a stated intention was a lie stays the surprisal's job, not the
classifier's.

Two limits come with reading talk this way. The five classes are fixed in advance
from standard cheap-talk categories rather than discovered in the
data, which keeps the distribution interpretable and lets it stand as a
pre-registered descriptive outcome, but means a message that fits none of them is
recorded as
neutral by default. And the labels come from a single model with no second rater,
and were never scored against human coding of the real messages, so unlike the
deception signal they carry no validated accuracy of their own. We therefore read
the class distribution as a descriptive instrument whose weight comes from
agreeing with the action-based measures, not from a per-message accuracy we never
established. The per-tier distribution is reported in Chapter~\ref{ch:results}.

\subsection{Behavioral tier inference}
\label{subsec:tier}

Where the message classifier of Section~\ref{subsec:classification} reads only an
agent's words, this metric reads only its play: it asks whether an agent's
capability tier can be recovered from how it acts, with the messages set aside. We
fit a logistic-regression classifier on per-agent behavioral features from the
Public Goods Game and read its predicted probability
$P(\mathrm{strong}\mid\mathrm{behavior})$ that an agent is strong. The target is
each agent's assigned capability tier, strong or weak, known by construction
rather than read off any outcome. The pre-registration lists this metric as a
secondary outcome under the name Bayesian tier inference; the label reflects the
posterior-style reading of the predicted probability rather than a model with
priors, and the implementation is the logistic regression described
here. It is scored by balanced accuracy and the area under the ROC curve, both
with a chance baseline of $0.5$ rather than the majority class.

The features are five summaries of an agent's contributions across the ten
rounds, all drawn from behavior and none from payoff, so that tier cannot leak in
through the outcome the classifier is meant to predict:
\begin{description}
  \item[mean contribution:] the average amount the agent puts into the pool;
  \item[contribution volatility:] the standard deviation of that amount across
    rounds, steady versus erratic play;
  \item[trend:] the slope of contribution against round number, drift up or down
    over the game;
  \item[end-game drop:] how far the final-round move falls below the agent's
    earlier average, the last-round defection spike;
  \item[reciprocity:] the correlation between the agent's move and what others
    contributed the round before.
\end{description}

Each agent supplies one feature row, so the sample is small and a single
train-test split would be noisy. We evaluate instead by stratified $k$-fold
cross-validation: the classifier is trained on a subset of agents and scored on
the held-out remainder, the folds rotate so every agent is scored once, and the
feature standardization is fit inside each fold so no test agent shapes its own
scaling. Balanced accuracy and the area under the ROC curve are averaged across
the folds; a score near that baseline would mean behavior carries no tier signal
at all.

That score carries two caveats. First, the evaluation is run on the
diverse strong tier by design. In the pilot the strong tier was a single model, so
a classifier could separate the tiers by fingerprinting that one model, and
accuracy was inflated. Running it on three rotated strong models changes the
question from telling one model from another to telling the strong tier from the
weak one, where a less-than-perfect score is the honest outcome: a classifier that
still reached ceiling would more likely have found a model signature than a tier. Second, the behavioral features here overlap
with those behind the interaction network of Section~\ref{subsec:network}, so the
two are not independent confirmations of tier separability; they are two views of
the same behavioral gap. The ROC curve and the production scores are reported in
Chapter~\ref{ch:results}.

\subsection{Interaction-network structure}
\label{subsec:network}

Where the tier classifier of Section~\ref{subsec:tier} turns each agent's play
into a feature vector, this last metric turns the agents into a graph and asks
about structure. The two games give two different graphs, and only one is a
genuine interaction. The Trust Game has literal transfers, an investor sending
tokens and a trustee returning them, so its graph is a directed
\emph{token-transfer network}: the nodes are agents and the edges carry the tokens
that actually moved between them. The Public Goods Game has no such transfers,
since contributions go into a shared pool rather than from one agent to another.
There we build a \emph{behavioral-similarity graph} instead, linking agents whose
contribution patterns resemble each other, and the pre-registration labels it as
similarity rather than interaction to keep the difference explicit.

On each graph we compute standard structural measures: centrality (degree,
betweenness, and eigenvector) for an agent's structural position, and tier
assortativity for whether agents of the same tier attach to one another. On the Trust transfers the reading that
matters is the asymmetry, which way the tokens net out between the strong and weak
tiers; on the similarity graph, assortativity measures how sharply the two tiers
pull apart.

We treat the network analysis as secondary. The Public Goods similarity graph is
not an interaction graph: an
edge means two agents played alike, not that anything passed between them, so it
cannot be read as who exploited whom, and only the Trust transfers carry that
meaning. The same graph is built from the behavioral features of the tier
classifier of Section~\ref{subsec:tier}, so the tier separation it shows, the
assortativity included, is that classifier's signal through a different lens, not
an independent confirmation. The graphs are also small and often fragmented: a
Trust game is two nodes and a Public Goods game four, and the sparsified
similarity graph breaks into disconnected pieces at scale, where eigenvector
centrality becomes ill-defined across the components and unreliable. We therefore
read the similarity graph through its tier assortativity rather than through
centrality, and the Trust graph through its transfer asymmetry, which needs no
centrality at all. The token-transfer figure and the tier assortativity are
reported in Chapter~\ref{ch:results}.

\section{Sample, seeds and budget}
\label{sec:sample}

The corpus is 320 games: the sixteen cells of Table~\ref{tab:matrix}, each run over
twenty random seeds. A third game, multi-issue negotiation, was planned and then
dropped after a feasibility probe found that the weak model could not reliably
value the offers, so the corpus spans the two games of
Section~\ref{sec:protocols} rather than three. Within a cell the twenty seeds are
paired across communication: a seed fixes the agents and their random draws, and
the communication flag is the only thing that changes between its on and off
games. That pairing makes the communication contrast, which the headline result
turns on, a within-seed comparison rather than a between-seed one. Because each game is fixed entirely by its configuration and one integer seed, the run
reproduces exactly, down to the anonymous agent identifiers and the shuffled agent
order in each prompt, so that neither an agent's name nor its position carries its
tier.

A hard cap governed the spending. A single tracker keeps the cumulative API cost
in the database and checks it before every paid call, raising an error that aborts
before any call would carry the total past 100 euros; the limit is never raised in
code, and reading the running total from the database means the cap holds across a
restart. It was never close. The 320 games cost 1.17 euros, and the paid
deception-scoring and message-classification passes brought the project to 5.62
euros in total, well under both the projection and the cap.

\section{Pre-registration and analysis plan}
\label{sec:prereg}

Everything in this chapter was fixed in advance. Before the first production game
ran, the design, the hypotheses, the analysis, and the criteria that would
falsify each claim were written up and registered on the Open Science Framework,
timestamped and locked against later editing (osf.io/rnqgs, 18~June 2026). Locking
does one narrow thing: a prediction registered before the data
exists and confirmed afterward is a confirmation, while the same prediction told
after the fact is only a story.

Five hypotheses were registered, in the numbering the results chapter keeps:
\begin{description}
  \item[H1] mixed-capability populations show higher welfare inequality than
    homogeneous ones;
  \item[H2] the strong tier does not systematically out-earn the weak tier, and
    may earn less, a negative capability premium;
  \item[H3] communication worsens the strong tier's position, a more negative
    premium with it than without;
  \item[H4] contemporaneous surprisal tracks human-judged deception better than
    the global policy-KL divergence, registered as exploratory;
  \item[H5] the direction of these effects holds across the two games rather than
    reversing between them.
\end{description}
Of the five hypotheses, H4 is registered as exploratory and the other four as
confirmatory predictions; the tier inference, message classification, and network
structure of Section~\ref{sec:metrics} are secondary, descriptive outcomes.

The analysis was pre-specified as well. The primary model is a mixed-effects
regression of payoff on composition, game, and communication, with random
intercepts for agent and for game to absorb the repeated agents and repeated
games. Inequality is read through the indices of
Section~\ref{subsec:inequality} with bootstrap confidence intervals, and
comparisons across the cell matrix are corrected for multiplicity by the
Benjamini--Hochberg procedure. The deception test is the one place the plan set a
hard bar: the contemporaneous surprisal had to reach a Spearman correlation of
$0.40$ with the human labels to count as confirmatory, a threshold fixed before
the pilot and held when the pilot fell below it, which is why surprisal is
reported as exploratory and policy-KL as a negative control
(Section~\ref{subsec:deception}). A tier-level
effect is trusted only when it holds across the three distinct strong models
rather than within one, which is what the rotation of Section~\ref{sec:tiers} was
built to allow.

\section{Integrity and design limitations}
\label{sec:integrity}

The pre-registration also recorded what the design cannot show, before the data
could tempt anyone to forget. The deception metric was validated by a single human
rater, who is also its author; a blind, frozen rubric and single-pass coding guard
against the obvious bias, but there is no second rater and so no inter-rater
reliability. The study is pilot-scale. The strong tier is three specific models,
so a claim about strong models in general is only ever a claim about these three.
And the deception finding is exploratory by registration, not confirmatory
(Section~\ref{subsec:deception}).

One limit was not in the pre-registration and is worth stating plainly. The model
that scores deception and classifies messages, Claude Haiku~4.5, is also one of
the three strong players, and it scores the games it played in along with the
others, so a reader could fairly call the arrangement circular, a competitor
judging the contest. But the surprisal and the policy-KL control are computed from
the same Haiku-inferred stated policy and differ only in what they compare it
against, the one a single round's action, the other the whole game's. A bias in
how the model infers that policy would land on both measures alike, so it cannot
be why the contemporaneous signal tracks deception and the global one does not;
that difference is built into their baselines, not into the reading. The blind
classifier of Section~\ref{subsec:classification} adds a check from outside that
machinery, never seeing the action behind a message, so it corroborates the
surprisal from text alone. The finding rests on that agreement, not on a single
read (Section~\ref{subsec:deception}).

A larger qualification concerns the communication condition, where the prompt and
the runner part ways. With communication on, each agent is asked to write a public
message and told the others will see it, but the runner never delivers it: an agent
sees the others' past actions, not their messages (Section~\ref{sec:communication}).
The condition is therefore public commitment without reception, and whatever it
changes in behavior it changes through the framing of being asked to pledge, not
through any exchange of cheap talk; the results and discussion read the effect on
that understanding. The messages keep their value for the deception measures, which
compare an agent's own words with its own deeds and so do not depend on another
player reading them.

Other bounds are argued where the metrics are defined and only gathered here.
Inequality indices on a two-player Trust game are coarse, so there the capability
premium carries the comparison (Section~\ref{subsec:inequality}). The network
analysis is secondary: its Public Goods graph is behavioral similarity rather than
interaction, and its tier signal is not independent of the behavioral classifier
(Sections~\ref{subsec:network} and~\ref{subsec:tier}). The message classifier's
five labels were never checked against human coding, so they are read as
description rather than measurement (Section~\ref{subsec:classification}). All of
these limits are recorded, in the pre-registration or here, not left for the
reader to find.
